Yao et al.~\cite{yao2023react} introduce ReAct, which interleaves reasoning
traces (\emph{Thought}) with actions (\emph{Search}, \emph{Lookup},
\emph{Finish}) against a Wikipedia-backed environment. Section 3.3 of the
paper additionally proposes two hybrid strategies that combine ReAct with
chain-of-thought self-consistency (CoT-SC): \emph{ReAct$\to$CoT-SC} backs off
from ReAct to a majority vote over $n$ independent CoT samples, and
\emph{CoT-SC$\to$ReAct} does the reverse. Both directions require a
\emph{trigger}: a rule for deciding when the first strategy has failed and
the second should be invoked instead.

For ReAct$\to$CoT-SC, the paper's trigger is a single condition: ReAct
exhausted its step budget without ever emitting a \texttt{Finish[\ldots]}
action. This is a reasonable proxy for one failure mode --- the agent
looping without converging --- but it says nothing about episodes in which
ReAct \emph{does} emit \texttt{Finish}, just on the basis of too little or
the wrong evidence.

This paper reports on a modification of that single entry point. Building
on a reduced-scale reproduction of the full ReAct pipeline (Assignment 2 of
this course), we show with trajectory-level evidence that the paper's
trigger is domain-dependent in a way that matters: it detects failures on
one of our two evaluation domains and is essentially inert on the other.
We generalize the trigger into a small set of independently-validated
signals, ablate them against an oracle correctness label \emph{before}
spending compute on a full comparison, and report both the positive and
negative results that came out of that validation step.

Our contributions are: (1) a quantified account, at $n=100$ per domain, of
\emph{why} the paper's back-off trigger under-fires on FEVER-style
claim verification; (2) three back-off signals evaluated independently
for precision lift over the domain's base error rate, showing that
evidence-count alone is not a useful proxy for evidence quality despite
being the cheaper signal to compute; (3) an entailment-verification signal
that is a useful proxy, at the cost of one additional LLM call per
question; and (4) an accuracy-vs-cost analysis of the resulting hybrid,
including a corrected token-truncation bug in the underlying reproduction
that on its own changed FEVER accuracy by a large margin, reported
separately from the trigger improvement's own credit.
